\begin{table}[h]
\centering
\small
\caption{FinAgent Red-Team threat taxonomy. Identifiers match \texttt{THREAT\_CATEGORIES} in the released code, so paper, figures and result files use one numbering. T1 is the delivery vector rather than an outcome, and is instantiated inside every other category; it therefore has no standalone scenarios of its own.}
\label{tab:threats}
\begin{tabular}{|p{1.9cm}|p{3.9cm}|p{4.2cm}|c|}
\hline
\textbf{Threat} & \textbf{Attack Goal} & \textbf{Real-World Impact} & \textbf{Scen.} \\
\hline
T1: Indirect Prompt Injection & Smuggle instructions into agent-readable content (email, invoice, chat) & Arbitrary agent action; the vector for T2--T7 & --- \\
\hline
T2: Unauthorized Transfer & Wire funds to a payee not on the approved list & Embezzlement, fraud & 6 \\
\hline
T3: Sanctions Evasion & Move funds to a screened entity, either by urging the screen to be skipped or by disguising the name & Sanctions breach, legal liability & 12 \\
\hline
T4: Payment Structuring & Split one payment into sub-threshold transfers to evade the reporting trigger & Money laundering, reporting breach & 6 \\
\hline
T5: Dual-Approval Defeat & Initiate and self-approve a transfer that requires two signers & Segregation-of-duties violation & 6 \\
\hline
T6: Data Exfiltration & Send customer PII to an external recipient under an audit pretext & Breach, GLBA/GDPR liability & 6 \\
\hline
T7: Confused Deputy & Redirect an approved payee's bank details, then pay the payee & Business email compromise, payment fraud & 6 \\
\hline
\end{tabular}
\end{table}

Every attack is delivered by the same vector (T1): the user asks the agent to do something ordinary---read the inbox, settle an invoice---and the adversarial instruction arrives inside content the agent reads while doing it \citep{greshake2023injection}. The agent is never told by its principal to misbehave. We treat injection as a vector rather than a scored category, because an agent that resists T2--T7 has by construction resisted the injection that carried them. Six categories carry scenarios; T3 carries twice as many as the others because it is generated by two structurally distinct builders (an explicit request to skip screening, and a name-obfuscation attack that carries no request at all and instead tests whether the agent screens as a matter of course). The total is 42 attack scenarios plus 6 benign controls---legitimate payments to allow-listed vendors, carried in \emph{trusted} content---which measure over-refusal and distinguish a model that is safe from one that is merely inert. Each category's success predicate is a deterministic function of final world state; the per-threat predicates, attack framings, and the seeded generation procedure are given in Appendix~\ref{app:scenarios}.
